% Gerado a partir de docs/tcc.md por scripts/sincronizar_tcc.py.
\chapter{Referencial Teórico}\label{referencial-teuxf3rico}

Avaliar um serviço público digital exige olhar para mais do que a tela final de confirmação. A experiência do cidadão começa quando ele tenta entender onde está o serviço, continua nos caminhos que percorre e termina apenas quando a demanda é resolvida ou abandonada. Por isso, este capítulo reúne os fundamentos que sustentam a proposta do avaliador: conceitos de serviço público, satisfação, qualidade percebida, usabilidade, modelos brasileiros de avaliação e sinais de possível dificuldade observáveis durante a navegação.

O capítulo também parte de uma tradição anterior à digitalização dos serviços: os modelos clássicos de qualidade em serviços. Essa base é importante porque mostra que a percepção do usuário sempre dependeu da comparação entre o que ele esperava receber e o que efetivamente experimentou. No ambiente público digital, essa comparação passa a envolver, além do resultado do serviço, a clareza da informação, a confiança no canal, a facilidade de navegação, o esforço exigido e a capacidade do órgão de responder ao problema identificado \cite{gronroos1984service,parasuraman1988servqual,deloneMclean2003}.

\section{Serviços e Serviços Públicos}\label{serviuxe7os-e-serviuxe7os-puxfablicos}

Serviços diferenciam-se de produtos porque são mais intangíveis, ocorrem como séries de atividades, envolvem algum grau de interação com o usuário e, em muitos casos, são produzidos e consumidos simultaneamente. Essa característica torna a avaliação de serviços mais dependente da experiência vivida pelo usuário, e não apenas de propriedades observáveis de um bem material \cite{gronroos1984service,limaPereira2020modulo1}.

No setor público, o serviço é parte do relacionamento entre cidadão e Estado. Ele pode envolver uma entrega individualizada, como a emissão de um documento, ou uma entrega de interesse coletivo, como segurança, educação, saúde e infraestrutura. Por isso, a avaliação de serviços públicos possui uma dimensão adicional: além de medir satisfação ou qualidade, ela se conecta ao exercício da cidadania, à prestação de contas e ao aperfeiçoamento da administração pública \cite{brasil2017lei13460,brasil2021lei14129,limaPereira2020modulo1}.

Com a digitalização, muitos serviços passaram a ser acessados por portais, aplicativos e sistemas integrados. Essa mudança amplia o alcance da prestação pública, mas também cria novas exigências de usabilidade, clareza, confiabilidade e disponibilidade. A usabilidade pode ser entendida como resultado de uso associado à efetividade, eficiência e satisfação em um contexto definido \cite{iso9241112018}. O usuário não avalia apenas o resultado final do serviço; ele também avalia o caminho percorrido, a linguagem utilizada, o esforço exigido, a facilidade de navegação e a capacidade do canal digital de resolver sua demanda.

\subsection{Serviço digital, interação e resultado pretendido}\label{serviuxe7o-digital-interauxe7uxe3o-e-resultado-pretendido}

Neste trabalho, serviço público digital é o arranjo de informações, procedimentos, pessoas, regras e recursos tecnológicos pelo qual o cidadão busca um resultado de interesse público. O portal é um canal desse arranjo. Encontrar uma página ou concluir uma interação não significa necessariamente obter o documento, benefício ou solução pretendida. A orientação institucional de resolver o problema completo do usuário ajuda a delimitar essa diferença \cite{govuk2026wholeProblem}.

A avaliação distingue quatro objetos: \textbf{qualidade do software}, relativa às propriedades da aplicação e da API; \textbf{qualidade da interação}, relativa às condições e resultados do uso da interface; \textbf{qualidade percebida do serviço}, relativa ao julgamento declarado sobre aspectos da prestação; e \textbf{efetividade do serviço}, relativa ao alcance do resultado pretendido. A primeira é tratada pelo modelo de qualidade de produto ISO/IEC 25010:2023; a usabilidade em contexto considera efetividade, eficiência e satisfação conforme ISO 9241-11; os atributos de serviço e a conclusão real exigem evidências próprias \cite{iso250102023,iso9241112018}.

Por exemplo, uma API pode aceitar e armazenar corretamente uma avaliação sem que o cidadão tenha conseguido declarar o imposto. Uma interface pode ser fácil de navegar, mas conduzir à consulta quando a necessidade era declarar. Uma nota favorável não prova que a obrigação foi cumprida. A conclusão de tarefa registrada no Figma é efetividade da tarefa simulada, não entrega de um serviço público real. O avaliador apoia investigação sobre a jornada; não executa nem certifica essa entrega.

\section{Avaliação de Serviços}\label{avaliauxe7uxe3o-de-serviuxe7os}

A avaliação de serviços pode ser entendida como uma investigação qualitativa ou quantitativa sobre a qualidade do serviço e a satisfação do usuário, considerando a relação entre expectativa e experiência concreta \cite{limaPereira2020modulo1}. Em outras palavras, avaliar um serviço significa observar se a entrega realizada corresponde ao que o usuário precisava, esperava ou considerava adequado. Modelos clássicos de qualidade de serviço indicam que a percepção do usuário é formada por lacunas entre o serviço esperado e o serviço recebido. O modelo de qualidade percebida de \citeonline{gronroos1984service} e o modelo SERVQUAL de \citeonline{parasuraman1988servqual} consolidaram a ideia de que qualidade em serviços depende de fatores como confiabilidade, responsividade, segurança, empatia e aspectos tangíveis. Embora esses modelos tenham origem no setor privado, eles influenciaram diversas iniciativas posteriores de avaliação de serviços públicos e serviços digitais.

No SERVQUAL, essas dimensões aparecem como tangibilidade, confiabilidade, responsividade, garantia e empatia. A contribuição central do modelo é tratar a qualidade percebida como diferença entre expectativa prévia e percepção após a experiência. Para serviços públicos, entretanto, essa leitura exige adaptação: o cidadão geralmente não escolhe entre prestadores concorrentes e muitas vezes depende daquele serviço para exercer um direito, cumprir uma obrigação ou acessar um benefício. Assim, a expectativa não é formada apenas por comparação de mercado, mas também por experiências anteriores com o Estado, comunicação institucional, relatos de outros usuários e, mais recentemente, pela comparação com serviços digitais privados que oferecem interfaces mais rápidas e responsivas \cite{parasuraman1988servqual,limaPereira2020modulo2}.

Nos serviços digitais, a avaliação passou a incorporar fatores ligados ao uso de sistemas de informação, como qualidade do sistema, qualidade da informação, qualidade do serviço, intenção de uso, satisfação do usuário e benefícios percebidos. Os materiais da Enap destacam que, à medida que cidadãos passaram a interagir com sistemas governamentais digitais, tornou-se necessário adotar modelos capazes de avaliar não apenas páginas ou portais, mas a experiência do usuário ao longo da prestação do serviço \cite{limaPereira2020modulo2,deloneMclean2003,pedrosa2019qualidadeGovbr}.

\section{Satisfação e Qualidade Percebida}\label{satisfauxe7uxe3o-e-qualidade-percebida}

Satisfação e qualidade percebida são conceitos próximos, mas não idênticos. A satisfação costuma expressar um julgamento global do usuário após consumir o serviço, enquanto a qualidade percebida pode ser decomposta em atributos mais específicos da experiência, como rapidez, clareza das informações, atendimento, esforço, confiabilidade e facilidade de uso. Para \citeonline{oliver2014satisfaction}, a satisfação está relacionada ao julgamento do usuário sobre o grau em que uma experiência atendeu ou não suas expectativas.

No contexto público, a distinção entre satisfação e qualidade percebida é importante porque uma nota geral, sozinha, nem sempre indica o que deve ser melhorado. Um usuário pode avaliar mal um serviço por demora, linguagem confusa, instabilidade do sistema, falta de atendimento ou excesso de etapas. Portanto, um avaliador de serviços públicos precisa combinar uma avaliação simples, como nota ou estrelas, com mecanismos que ajudem a qualificar o motivo da avaliação.

Essa lógica aparece na evolução dos instrumentos brasileiros de avaliação. O BRASP considera a qualidade percebida como um elemento central para compreender a experiência do usuário e apoiar o aprimoramento dos serviços \cite{silva2019brasp}. A proposta de modelo unificado busca integrar qualidade técnica, qualidade percebida, efetividade e satisfação, reduzindo o custo de resposta para o cidadão e tornando os resultados mais úteis para a gestão \cite{silva2022modeloUnificado}.

\section{Qualidade de Serviços Eletrônicos e Governo Eletrônico: E-S-QUAL, E-RecS-QUAL e e-GovQual}\label{qualidade-de-serviuxe7os-eletruxf4nicos-e-governo-eletruxf4nico-e-s-qual-e-recs-qual-e-e-govqual}

A passagem dos modelos gerais para os serviços eletrônicos exige considerar a prestação mediada por sistemas. \citeonline{parasuraman2005esqual} desenvolveram E-S-QUAL para sites de compras, com 22 itens e quatro dimensões: eficiência, cumprimento da prestação, disponibilidade do sistema e privacidade. E-RecS-QUAL trata situações não rotineiras e recuperação do serviço, com 11 itens em responsividade, compensação e contato. São escalas distintas; uma avaliação de recuperação pressupõe que tenha ocorrido uma situação pertinente.

A aplicação ao presente artefato é uma decisão de análise, não a aplicação integral dessas escalas. Eficiência motiva examinar caminhos e facilidade declarada; disponibilidade motiva testes técnicos; cumprimento motiva separar navegação de resultado; privacidade motiva requisitos de coleta. Na recuperação, responsividade e contato orientam investigar suporte e retorno ao cidadão. Compensação depende da natureza do serviço e de seu enquadramento institucional; não é inferida dos cliques ou de um comentário.

Os seis atributos do formulário não reproduzem os itens de E-S-QUAL/E-RecS-QUAL nem geram seus escores. O CES é outro indicador, e não uma dimensão substituta dessas escalas. Para serviços públicos, a interpretação precisa considerar direitos, obrigações, continuidade entre canais e ausência de escolha de prestador. Essa delimitação conduz aos modelos brasileiros apresentados nas seções seguintes, preservando a progressão qualidade de serviços \ensuremath{\rightarrow} serviços eletrônicos \ensuremath{\rightarrow} serviços públicos digitais \ensuremath{\rightarrow} BRASP e modelos nacionais.

No domínio governamental, \citeonline{papadomichelaki2012egovqual} desenvolveram e-GovQual para sites públicos nos quais o cidadão procura informação ou serviço. A escala final reúne 21 itens em confiabilidade, eficiência, suporte ao cidadão e confiança. Esse modelo oferece o elo entre qualidade eletrônica e prestação pública digital, sem reduzir a avaliação à transação comercial.

Neste trabalho, e-GovQual orienta examinar a consistência da prestação, a facilidade de acesso à informação, os caminhos de suporte e a confiança na coleta. São relações conceituais de projeto: o formulário não reproduz seus 21 itens, e os seis atributos não geram um escore e-GovQual. A validação relatada por Papadomichelaki e Mentzas não se transfere ao CES adaptado nem a este protótipo. O modelo internacional antecede, nesta progressão, os instrumentos brasileiros e o BRASP discutidos a seguir.

\section{Avaliação de Serviços Públicos no Brasil}\label{avaliauxe7uxe3o-de-serviuxe7os-puxfablicos-no-brasil}

A avaliação de serviços públicos no Brasil possui base normativa e base técnica. Do ponto de vista normativo, a Lei n\textsuperscript{o} 13.460/2017 estabelece direitos dos usuários e prevê que os resultados de avaliação subsidiem a reorientação dos serviços \cite{brasil2017lei13460}. O Decreto n\textsuperscript{o} 9.094/2017 institui a Carta de Serviços ao Usuário e reforça o uso de ferramenta de pesquisa de satisfação no governo federal \cite{brasil2017decreto9094}. A Lei n\textsuperscript{o} 14.129/2021 relaciona o governo digital à eficiência pública, à desburocratização e à participação do cidadão \cite{brasil2021lei14129}.

A partir desse conjunto normativo, a avaliação deixa de ser apenas uma prática gerencial desejável e passa a compor o próprio modo de prestação dos serviços públicos digitais. A Lei de Acesso à Informação reforça a transparência como base para acompanhamento social; a Lei n\textsuperscript{o} 13.460/2017 vincula avaliação à proteção do usuário; a Lei do Governo Digital incorpora canais digitais de participação; e normas infralegais detalham instrumentos de simplificação, monitoramento e avaliação de satisfação. O Quadro \ref{tab:marcos-legais-avaliacao} resume os marcos mais relevantes para este trabalho.

\begingroup\footnotesize\SingleSpacing
\setlength{\tabcolsep}{4pt}
\setlength{\LTpre}{.6\baselineskip}
\setlength{\LTpost}{.4\baselineskip}
\setlength{\LTcapwidth}{\textwidth}

\begin{TCCquadro}
\begin{longtable}[]{@{}
  >{\RaggedRight\arraybackslash}p{(\linewidth - 4\tabcolsep) * \real{0.2800}}
  >{\RaggedRight\arraybackslash}p{(\linewidth - 4\tabcolsep) * \real{0.1200}}
  >{\RaggedRight\arraybackslash}p{(\linewidth - 4\tabcolsep) * \real{0.6000}}@{}}
\caption[]{\TCCquadroentry{Marcos legais nacionais sobre avaliação de serviços públicos digitais}Marcos legais nacionais sobre avaliação de serviços públicos digitais}\label{tab:marcos-legais-avaliacao}\tabularnewline
\toprule\noalign{}
\begin{minipage}[b]{\linewidth}\raggedright
\textbf{Instrumento}
\end{minipage} & \begin{minipage}[b]{\linewidth}\raggedright
\textbf{Ano}
\end{minipage} & \begin{minipage}[b]{\linewidth}\raggedright
\textbf{Principais disposições relevantes}
\end{minipage} \\
\midrule\noalign{}
\endfirsthead
\caption[]{Marcos legais nacionais sobre avaliação de serviços públicos digitais (continuação)}\tabularnewline
\toprule\noalign{}
\begin{minipage}[b]{\linewidth}\raggedright
\textbf{Instrumento}
\end{minipage} & \begin{minipage}[b]{\linewidth}\raggedright
\textbf{Ano}
\end{minipage} & \begin{minipage}[b]{\linewidth}\raggedright
\textbf{Principais disposições relevantes}
\end{minipage} \\
\midrule\noalign{}
\endhead
\bottomrule\noalign{}
\endlastfoot
Lei n\textsuperscript{o} 12.527 (LAI) \cite{brasil2011lei12527} & 2011 & Transparência ativa e passiva; base para acesso a informações públicas e dados sobre desempenho de serviços. \\
Lei n\textsuperscript{o} 13.460 \cite{brasil2017lei13460} & 2017 & Direitos do usuário, participação na avaliação dos serviços e uso dos resultados para reorientação da prestação pública. \\
Decreto n\textsuperscript{o} 9.094 \cite{brasil2017decreto9094} & 2017 & Simplificação do atendimento, Carta de Serviços ao Usuário e racionalização de exigências feitas ao cidadão. \\
Lei n\textsuperscript{o} 14.129 \cite{brasil2021lei14129} & 2021 & Governo digital, eficiência pública, desburocratização, canais digitais de participação e melhoria da prestação estatal. \\
Decreto n\textsuperscript{o} 11.260 \cite{brasil2022decreto11260} & 2022 & Referência histórica sobre elaboração da estratégia e prorrogação da anterior; revogado em 2024, sem caráter de norma vigente neste trabalho. \\
Portaria SGD/ME n\textsuperscript{o} 548 \cite{brasil2022portaria548} & 2022 & Avaliação de satisfação dos usuários e padrões de qualidade para serviços públicos digitais. \\
\end{longtable}
\end{TCCquadro}

\endgroup

Fonte: elaborada pelos autores com base na legislação federal brasileira.

Do ponto de vista técnico, relatórios produzidos no âmbito da Universidade de Brasília, da Secretaria de Governo Digital e de instituições parceiras apresentam modelos voltados à mensuração da qualidade de serviços públicos digitais. A proposta Qualidade Gov.br organiza a avaliação em dimensões ligadas à consistência digital e à experiência do usuário, reconhecendo que a qualidade de um serviço digital depende tanto dos requisitos técnicos quanto da jornada do cidadão \cite{pedrosa2019qualidadeGovbr}. O BRASP amplia essa discussão ao considerar serviços presenciais, digitizados e parcialmente digitizados \cite{silva2019brasp}.

O avanço desses modelos demonstra que a avaliação pública precisa equilibrar rigor metodológico e simplicidade operacional. Instrumentos longos e complexos podem gerar dados ricos, mas tendem a aumentar o esforço do respondente. Por outro lado, instrumentos muito curtos podem ampliar a adesão, mas produzir dados insuficientes para orientar melhorias. Esse equilíbrio é especialmente relevante para a proposta deste TCC, que busca um avaliador simples para o cidadão e útil para a gestão.

\section{BRASP, API de Avaliação e Modelo Unificado}\label{brasp-api-de-avaliauxe7uxe3o-e-modelo-unificado}

O Modelo Brasileiro de Avaliação de Serviços Públicos (BRASP) foi desenvolvido para construir e validar instrumentos de avaliação da qualidade adaptados à realidade brasileira. O modelo considera três modalidades de prestação: serviços digitizados, serviços parcialmente digitizados e serviços presenciais \cite{silva2019brasp}. Essa diferenciação é relevante porque a experiência do usuário muda de acordo com o canal. Um serviço presencial pode depender mais de atendimento, espera e estrutura física; um serviço digital pode depender mais de facilidade de uso, clareza, disponibilidade e segurança; e um serviço híbrido pode combinar problemas de ambos os contextos.

Nos materiais da Enap, o modelo brasileiro de avaliação é apresentado como instrumento para identificar fatores relacionados à satisfação do usuário, localizar pontos críticos e gerar comparações capazes de acompanhar a evolução dos serviços. Essa leitura é coerente com o BRASP e com a proposta de modelo unificado, que organizam satisfação, qualidade percebida, qualidade técnica e efetividade como dimensões complementares de análise \cite{limaPereira2020modulo3,silva2019brasp,silva2022modeloUnificado}. A pesquisa de satisfação aparece ancorada em uma pergunta geral sobre a experiência com o serviço, enquanto a qualidade percebida é investigada por atributos que ajudam a explicar essa experiência.

A API de avaliação do governo brasileiro segue uma lógica semelhante. Segundo o relatório sobre os atributos da API, a satisfação é medida por uma estratégia simples inspirada no CSAT, com resposta por opções associadas a estrelas, e a qualidade percebida é explorada por atributos como clareza das informações, custo ou esforço para obtenção do serviço, disponibilidade, eficácia, facilidade de uso, qualidade dos canais de comunicação e tempo para obtenção do serviço \cite{silva2024atributosApi}. Essa abordagem se aproxima da proposta deste trabalho, pois combina avaliação rápida com categorização do motivo da nota.

A proposta de modelo unificado reforça essa direção ao articular quatro eixos: qualidade técnica, qualidade percebida, satisfação do usuário e efetividade \cite{silva2022modeloUnificado}. Para um avaliador de serviços públicos, esses eixos podem orientar tanto a coleta quanto a visualização de dados. A satisfação indica o julgamento geral do cidadão; a qualidade percebida explica a experiência por atributos; a qualidade técnica pode apoiar análises de disponibilidade, integração e consistência; e a efetividade aproxima a avaliação da pergunta central: o serviço resolveu a necessidade do usuário?

\section{Usabilidade, Inclusão Digital e Métricas de Experiência}\label{usabilidade-inclusuxe3o-digital-e-muxe9tricas-de-experiuxeancia}

A usabilidade em serviços públicos digitais possui impacto mais amplo do que em produtos privados. Quando uma interface privada é difícil, o usuário pode trocar de fornecedor; quando um serviço público digital é difícil, o cidadão pode perder prazo, deixar de acessar um benefício, não concluir uma obrigação ou depender de atendimento presencial. Por isso, os princípios de usabilidade precisam ser interpretados como parte do acesso efetivo ao serviço. As heurísticas de \citeonline{nielsen1994heuristic}, como visibilidade do status do sistema, correspondência entre linguagem da interface e linguagem do usuário, prevenção de erros e ajuda contextual, são diretamente relacionadas aos problemas que o avaliador busca observar durante a navegação.

A representatividade também é um ponto relevante. Serviços públicos atendem usuários com diferentes níveis de letramento digital, conectividade, idade, familiaridade com formulários e capacidade de compreender linguagem institucional. Os dados do PPA Participativo 2024-2027 indicam que participação digital em larga escala pode conviver com barreiras de uso e de retorno ao cidadão \cite{brasilMgi2024ppaParticipativo}. Da mesma forma, pesquisas sobre acesso às tecnologias de informação no Brasil mostram que o uso de serviços digitais ainda depende de condições desiguais de acesso, dispositivos e habilidades \cite{cetic2025ticDomicilios2024}. Para este TCC, isso reforça a necessidade de uma coleta de feedback simples, opcional e de baixo esforço.

O CSAT e o CES abordam julgamentos diferentes. A satisfação expressa a avaliação do cidadão sobre a experiência com a etapa do serviço; o CES aborda a facilidade ou o esforço percebido nessa interação. A proposta usa estrelas como operacionalização local de satisfação, atributos para detalhar aspectos do serviço e um item CES para facilidade percebida. Nenhum deles é obtido diretamente da telemetria. A discussão profissional de \citeonline{dixon2010effort} motiva atenção ao esforço, mas sua origem em atendimento comercial não demonstra validade em serviços públicos.

A literatura profissional de \citeonline{dixon2013effortless} desenvolve a abordagem de experiência com menor esforço no atendimento comercial. Para a formulação por concordância, a apresentação de \citeonline[slide 72]{dixon2016ces} mostra o CES 2.0 com afirmação sobre facilidade para lidar com uma questão. Esse material primário distingue a formulação posterior da referência de 2010, que não é utilizada para atribuir origem à redação ou à escala adotadas aqui. A conferência visual desse slide identifica seis opções visíveis, não sete; ele fundamenta a formulação por concordância, não o número de categorias adotado neste projeto. A frase em português, o recorte por etapa e as sete categorias são decisões de adaptação deste projeto, ainda sem validação com cidadãos.

O CES adotado usa a afirmação ``Foi fácil realizar o que eu precisava nesta etapa do serviço'', com sete categorias: 1 --- discordo totalmente; 2 --- discordo; 3 --- discordo parcialmente; 4 --- nem concordo nem discordo; 5 --- concordo parcialmente; 6 --- concordo; 7 --- concordo totalmente. A resposta é opcional. Notas maiores significam maior facilidade declarada e, nessa orientação, menor esforço percebido. Não se inverte o resultado nem se combina essa escala com medidas de outra orientação. Trata-se de adaptação ao contexto deste protótipo, não de uma tradução validada. A documentação de \citeonline{surveymonkey2026ces} é utilizada como apoio comercial complementar à descrição operacional, sem importar alegações comerciais de desempenho ou pontos de corte.

A média CES descritiva é a soma das respostas válidas dividida pelo número de respostas válidas, com apresentação conjunta da distribuição nas sete categorias e do denominador. Ausência de resposta fica nula, nunca zero; na ausência de respostas, a média fica indisponível. A média resume uma escala ordinal por convenção descritiva do protótipo e não comprova intervalos equivalentes entre categorias. Não se adotam faixas universais de ``bom'' ou ``ruim'' nem um escore geral resultante da soma de estrelas, atributos, CES e telemetria.

O item se refere à etapa vivenciada até o convite, inclusive quando a tarefa não foi concluída. A expressão ``realizar o que eu precisava'' pode, porém, ser interpretada como tarefa já concluída; o piloto de compreensão deverá verificar esse risco antes de manter a redação em convite intermediário. Portanto, facilidade declarada não comprova entrega do resultado pretendido pelo cidadão. O momento da coleta pode influenciar recordação e interpretação. Resultados sobre avaliação retrospectiva motivam investigar esse efeito, sem demonstrar a vantagem do acionamento neste cenário \cite{kahneman1993peakEnd}. O esforço para preencher o próprio formulário é outra variável e requer medida específica.

\section{Experiência do Usuário e Avaliações Inspiradas em Plataformas Digitais}\label{experiuxeancia-do-usuuxe1rio-e-avaliauxe7uxf5es-inspiradas-em-plataformas-digitais}

Plataformas digitais privadas popularizaram formas de avaliação de baixo esforço: nota de uma a cinco estrelas, comentários curtos, seleção de motivos e visualização agregada das avaliações. A referência a aplicações como plataformas de entrega deve ser entendida como inspiração de interação, e não como modelo de governança. Em uma plataforma privada, a avaliação frequentemente orienta reputação, consumo e concorrência. Em serviços públicos, ela deve orientar direitos, transparência, participação social e melhoria da entrega estatal.

Mesmo com essa diferença, a familiaridade do usuário com mecanismos simples de avaliação pode ser aproveitada. Um avaliador de serviços públicos pode reduzir barreiras de participação se oferecer uma interface direta, com poucos passos e linguagem compreensível. A nota geral permite captar a percepção imediata; os atributos permitem interpretar os motivos; e comentários opcionais podem revelar problemas não previstos pelo formulário. Essa escolha está alinhada ao design centrado no usuário, que orienta o desenvolvimento de sistemas interativos a partir das necessidades, capacidades e limitações das pessoas que irão utilizá-los \cite{iso92412102019}.

Para que a avaliação seja útil ao gestor, o sistema também precisa organizar os dados coletados. Isso envolve agrupar avaliações por serviço, período, órgão, canal e atributo; identificar itens recorrentes; apresentar médias e distribuições; e permitir comparação temporal. A literatura de melhoria de serviços reforça que as avaliações devem alimentar ciclos de melhoria contínua, funcionando como linha de base para decisões e como indicador para acompanhar a evolução do serviço \cite{limaPereira2020modulo4,govuk2021usingData,govbrFerramentaAvaliacao}.

\subsection{Posicionamento em relação a soluções existentes}\label{posicionamento-em-relauxe7uxe3o-a-soluuxe7uxf5es-existentes}

Embora o protótipo dialogue com ferramentas existentes de avaliação e análise de experiência, sua contribuição não está em substituí-las nem em afirmar superioridade. A contribuição está em combinar, em um MVP com regras explícitas e voltado a serviços públicos, sinais simples de possível dificuldade durante a navegação, coleta contextual por formulário curto e organização administrativa dos dados por serviço e atributo.

A ferramenta oficial de avaliação do gov.br constitui uma referência institucional para medir satisfação em serviços públicos digitais. Ferramentas comerciais como Hotjar e Microsoft Clarity oferecem funcionalidades de analytics comportamental, mapas de calor, gravações de sessão e feedback. Suas páginas de documentação são utilizadas neste trabalho somente para caracterização tecnológica das funcionalidades declaradas, e não como evidência científica de eficácia ou adequação ao setor público \cite{govbrFerramentaAvaliacao,hotjar2026,microsoftClarity2026}. O Government Digital Service, por sua vez, é uma fonte institucional e metodológica sobre uso de dados de desempenho na melhoria de serviços digitais \cite{govuk2021usingData}.

O Quadro \ref{tab:posicionamento-solucoes} sintetiza esse posicionamento.

\begingroup\footnotesize\SingleSpacing
\setlength{\tabcolsep}{4pt}
\setlength{\LTpre}{.6\baselineskip}
\setlength{\LTpost}{.4\baselineskip}
\setlength{\LTcapwidth}{\textwidth}

\begin{TCCquadro}
\begin{longtable}[]{@{}
  >{\RaggedRight\arraybackslash}p{(\linewidth - 6\tabcolsep) * \real{0.1800}}
  >{\RaggedRight\arraybackslash}p{(\linewidth - 6\tabcolsep) * \real{0.2300}}
  >{\RaggedRight\arraybackslash}p{(\linewidth - 6\tabcolsep) * \real{0.2500}}
  >{\RaggedRight\arraybackslash}p{(\linewidth - 6\tabcolsep) * \real{0.3400}}@{}}
\caption[]{\TCCquadroentry{Posicionamento da proposta em relação a soluções existentes}Posicionamento da proposta em relação a soluções existentes}\label{tab:posicionamento-solucoes}\tabularnewline
\toprule\noalign{}
\begin{minipage}[b]{\linewidth}\raggedright
\textbf{Solução}
\end{minipage} & \begin{minipage}[b]{\linewidth}\raggedright
\textbf{Foco principal}
\end{minipage} & \begin{minipage}[b]{\linewidth}\raggedright
\textbf{Contribuição}
\end{minipage} & \begin{minipage}[b]{\linewidth}\raggedright
\textbf{Limite em relação a este TCC}
\end{minipage} \\
\midrule\noalign{}
\endfirsthead
\caption[]{Posicionamento da proposta em relação a soluções existentes (continuação)}\tabularnewline
\toprule\noalign{}
\begin{minipage}[b]{\linewidth}\raggedright
\textbf{Solução}
\end{minipage} & \begin{minipage}[b]{\linewidth}\raggedright
\textbf{Foco principal}
\end{minipage} & \begin{minipage}[b]{\linewidth}\raggedright
\textbf{Contribuição}
\end{minipage} & \begin{minipage}[b]{\linewidth}\raggedright
\textbf{Limite em relação a este TCC}
\end{minipage} \\
\midrule\noalign{}
\endhead
\bottomrule\noalign{}
\endlastfoot
Ferramenta de avaliação do gov.br & Avaliação institucional de satisfação & Padronização oficial no governo digital brasileiro & Não enfatiza, no escopo analisado, gatilhos baseados em cliques e tempo durante a jornada. \\
Hotjar & Analytics comportamental, mapas de calor, gravações e feedback & Observação de comportamento em sites & Documentação comercial; requer adaptação à governança pública e à minimização de dados. \\
Microsoft Clarity & Análise de sessões, mapas de calor e navegação & Observabilidade detalhada de páginas web & Documentação comercial; não constitui modelo de avaliação de serviço público. \\
GDS UK & Medição e melhoria de serviços digitais & Referência metodológica institucional para uso de dados & Contexto institucional distinto e sem equivalência direta ao protótipo. \\
Artefato deste TCC & Implementação local Python e interfaces de avaliação no Figma & Implementação parcial de eventos, feedback e painel; observações exploratórias de interface & Inspeção interna da implementação e avaliação de interface Figma/\allowbreak Useberry são procedimentos distintos. O estudo com participantes não verifica a API nem valida o serviço em contexto real. \\
\end{longtable}
\end{TCCquadro}

\endgroup

Fonte: elaborada pelos autores (2026).

A comparação é documental e de posicionamento: baseia-se nas funcionalidades declaradas pelas fontes institucionais e comerciais, sem teste independente ou auditoria técnica das soluções. Assim, delimita a lacuna explorada pelo trabalho sem transformar documentação de ferramentas em evidência científica.

\section{Métricas de Navegação e Sinais de Dificuldade}\label{muxe9tricas-de-navegauxe7uxe3o-e-sinais-de-dificuldade}

Além do feedback declarado pelo usuário, serviços digitais podem ser avaliados por sinais observados durante a jornada. A taxa de conclusão, por exemplo, mede a proporção de transações iniciadas que chegam ao ponto definido como conclusão do serviço; sua leitura depende de identificar início, etapas intermediárias, abandono e término do fluxo \cite{govuk2021completionRate}. Métricas de tempo, cliques, páginas visitadas e retorno a etapas anteriores também podem indicar dificuldade quando comparadas a faixas esperadas.

O uso dessas métricas aproxima a avaliação de serviços públicos da avaliação de sistemas de informação. O modelo de sucesso de sistemas de informação de \citeonline{deloneMclean2003} relaciona qualidade do sistema, qualidade da informação, qualidade do serviço, uso, satisfação do usuário e benefícios líquidos. Para a solução proposta, cliques, tempo e sequência de navegação são evidências comportamentais de possível esforço ou dificuldade. Eles não medem qualidade percebida. Estrelas, atributos e comentários constituem evidências declaradas sobre a experiência. O CES é autorrelato de facilidade percebida, e não uma transformação da quantidade de cliques ou do tempo. A contribuição analítica está em relacionar essas fontes sem fundi-las: a telemetria localiza o contexto e o feedback permite interpretar a percepção.

O registro de eventos de navegação exige cuidado metodológico e jurídico. A coleta deve ser limitada ao necessário para a finalidade de melhoria do serviço, com uso de identificador técnico e minimização dos campos registrados, sem presumir anonimização. A Lei Geral de Proteção de Dados Pessoais estabelece princípios para o tratamento de dados, inclusive em meios digitais, com objetivo de proteger liberdade, privacidade e desenvolvimento da personalidade da pessoa natural \cite{brasil2018lgpd}. Por isso, a proposta deste TCC considera o registro de cliques e tempo como mecanismo de análise agregada, não como instrumento de vigilância individual.

\section{Feedback Contextual e Acionamento por Sinais de Uso}\label{sec:feedback-contextual}

O feedback contextual busca aproximar a coleta declarada do ponto da jornada em que uma experiência ocorre. Essa proposta dialoga com o Experience Sampling Method (ESM), desenvolvido para coletar autorrelatos durante a vida cotidiana \cite{larson1983experienceSampling,hektner2007experienceSampling}. No presente trabalho, essa relação é conceitual: o convite acionado por um evento de navegação não representa, por si só, a aplicação integral de um protocolo de ESM.

A proximidade temporal pode reduzir a dependência de recordações retrospectivas e permitir investigar relatos mais situados. Estudos sobre avaliação retrospectiva mostram que o desfecho pode influenciar o julgamento de uma experiência \cite{kahneman1993peakEnd}; esse resultado não comprova que perguntar durante uma dificuldade de navegação produza respostas melhores. A formulação, as opções de resposta e o contexto da pergunta também influenciam o autorrelato \cite{schwarz1999selfReports}.

A eventual vantagem de obter um comentário mais específico precisa ser confrontada com o custo da interrupção. O estudo de \citeonline{mark2008interruptedWork} identificou maior estresse associado a interrupções em tarefas de trabalho. Esse achado oferece um motivo para investigar a intrusão do convite, sem estabelecer seu efeito no cenário deste TCC. Por isso, a pergunta contextual deve ser curta e opcional, e sua adequação precisa ser avaliada pelos usuários.

No protótipo deste TCC, os gatilhos comportamentais foram reduzidos a regras simples e explícitas: excesso de cliques, tempo acima da faixa esperada e persistência do sinal após a sugestão de redirecionamento. Esses sinais não substituem a percepção do usuário, mas indicam quando faz sentido solicitar uma avaliação de baixo esforço. A escolha por regras explícitas preserva a rastreabilidade do MVP e evita que a solução dependa de modelos estatísticos complexos antes de haver dados reais suficientes para calibração.

Esse tipo de coleta também precisa respeitar a carga cognitiva do cidadão. Questionários longos ou tecnicamente ambíguos aumentam esforço, erro e abandono. A Teoria da Carga Cognitiva indica que a memória de trabalho humana possui capacidade limitada; portanto, interfaces de feedback devem usar linguagem direta, poucas perguntas, opções de resposta claras e campos textuais opcionais \cite{sweller1988cognitiveLoad}. O formulário atual reúne satisfação por estrelas, atributos, CES opcional e comentário opcional. Essa decisão de projeto não demonstra esforço reduzido ou utilidade dos dados, propriedades que dependem de avaliação específica.

\section{Síntese do Referencial}\label{suxedntese-do-referencial}

O referencial teórico indica que a proposta de um avaliador de serviços públicos deve considerar quatro princípios. O primeiro é a centralidade do cidadão: a avaliação precisa capturar a experiência real do usuário, e não apenas métricas internas da administração. O segundo é a decomposição da avaliação em atributos: uma nota geral é útil, mas precisa ser acompanhada de critérios que expliquem o motivo da satisfação ou insatisfação. O terceiro é a orientação para melhoria: os dados coletados devem ser estruturados para apoiar decisões de gestão e acompanhar a evolução dos serviços. O quarto é a contextualidade como hipótese de desenho: investigar se aproximar a pergunta de um sinal de possível dificuldade altera a participação e o conteúdo das respostas, sem presumir que essa aproximação produza evidências melhores.

Com base nesses princípios, este TCC propõe um avaliador inspirado em interfaces simples de plataformas digitais, mas fundamentado em modelos de avaliação de serviços públicos, usabilidade, design centrado no usuário e sucesso de sistemas de informação. A solução deve permitir que o usuário avalie um serviço de forma rápida e compreensível, enquanto eventos de navegação e resultados agregados oferecem ao gestor indícios sobre quais aspectos do serviço precisam de priorização e aperfeiçoamento.
